% Fragment intended to be \input into a larger document.
% Deliberately uses no packages and defines no macros: it relies only on
% sectioning, \texttt, \emph, itemize and verbatim, all of which are part of
% core LaTeX. Adjust the sectioning level to taste.

\section{Modifications to Self-Collaboration}

The first of the two code generation solutions under evaluation is
Self-Collaboration, a multi-agent method in which three roles cooperate on a
programming task: an \emph{Analyst}, who inspects the repository and decides
what has to be written; a \emph{Coder}, who writes it; and a \emph{Tester},
who runs the tests and reports failures back to the Coder for a further round. The implementation we
evaluate is a fork of the authors' code, and our aim is to leave the method
itself untouched, so that differences in measured performance are attributable
to the method and the model rather than to our harness. Four interventions
were nevertheless necessary before the method could run as described and be
compared across configurations. The first is a matter of configuration. The
remaining three are changes to the implementation, and are the reason this
section exists at all.

A common cause underlies most of them. The implementation was written for
repair tasks, where an existing repository is modified under version control,
and it uses the resulting diff as its evidence that anything happened. The
tasks in this work are the opposite: an empty workspace and a specification
from which an entire project must be written. Nothing the agent produces is
tracked, so every mechanism that reads a diff reads nothing, and each such
mechanism fails in a different way. What we changed is the harness's reading
of the repository, never the roles or what they are asked to do.

The exact revision of the solution used in every run is recorded with the run
itself, so that any reported result can be traced to the code that produced it.

\subsection{Enabling the Tester Role}

In its default configuration the implementation never executes its Tester role.
Two independent conditions in the session loop must both hold before the Tester
is reached:

\begin{itemize}
\item A test command must be supplied. Where the caller passes none, the
      session leaves the loop immediately after the Coder, and the Tester is
      skipped.
\item More than one Coder round must be permitted. The Tester runs
      \emph{between} Coder rounds, so a session limited to a single round
      terminates before the Tester is reached, whatever test command was given.
\end{itemize}

Neither condition is met by the defaults, with the consequence that the
published three-role method degrades silently to a two-role one: the Analyst
localises the work, the Coder attempts it once, and nothing verifies the
result. This is not a defect in the solution but an omission in how it is
invoked, and we correct it in our experiment configuration rather than in the
solution's code. Both quantities are exposed as configuration values, are
recorded with each run, and can therefore be varied as hyperparameters.

The command given to the Tester installs the generated package and runs the
test suite that the agent itself has written. It has no access to the
benchmark's reference tests, which remain inside a separate container and are
never visible to the agent; the Tester's purpose is self-verification, not
grading.

\subsection{Change Detection on Tasks Built From Scratch}

The second intervention is a one-line change to the implementation.

After running the tests, the Tester performs a sanity check: it inspects the
repository to determine whether the Coder actually produced anything, and if it
finds no change it discards the test outcome and substitutes a synthetic
failure carrying the message ``No code changes were made''. The purpose of the
check is to prevent a Coder that exhausted its step budget without editing
anything from being credited with a passing test run.

The check was implemented with \texttt{git diff}:

\begin{verbatim}
diff_check = subprocess.run(
    "git diff", shell=True, cwd=self.repo_path,
    capture_output=True, text=True,
)
has_diff = bool(diff_check.stdout.strip())
if not has_diff:
    status = "failed"
    output = "No code changes were made - nothing to test."
\end{verbatim}

This is exact for the benchmark the implementation was written against.
SWE-bench presents an existing repository whose files are already under version
control; the Coder modifies those tracked files, and the resulting diff is not
merely evidence of work but the deliverable itself, since a SWE-bench
prediction \emph{is} a patch, which the harness extracts with the same command.
An empty diff there genuinely means that nothing was done.

The benchmark used in this work has the opposite shape. Each task supplies a
natural-language specification and an empty workspace, and the agent must
create an entire library from scratch. Every file the Coder writes is therefore
untracked. Because \texttt{git diff} compares the working tree against the
index, and untracked files appear in neither, the command reports nothing at
all, however much code was generated. The consequences compound: a genuine test
result is discarded on every round; the Tester can never report success, since
the substituted status is always a failure; and the report handed to the next
Coder round asserts that its predecessor's work does not exist, while
instructing it not to repeat the change that failed. The Tester would run, and
would systematically misreport what it found.

It is worth noting that this cannot be repaired by having the Coder commit its
work, which is the intuitive remedy. \texttt{git diff} reports unstaged
modifications to tracked files; after a commit the working tree matches the
index and the diff is empty again, so instructing the Coder to commit would
cause the check to fail deterministically rather than intermittently. Staging
alone fares no better, for the same reason. Neither is the method's concern in
any case: version control appears nowhere in the published description of
Self-Collaboration, and nowhere in the prompts given to any of its three roles.
It enters the implementation solely as the harness's mechanism for extracting a
SWE-bench patch.

We therefore changed the check to use \texttt{git status --porcelain}, which
reports untracked files alongside modifications:

\begin{verbatim}
change_check = subprocess.run(
    "git status --porcelain", shell=True, cwd=self.repo_path,
    capture_output=True, text=True,
)
has_changes = bool(change_check.stdout.strip())
if not has_changes:
    status = "failed"
    output = "No code changes were made - nothing to test."
\end{verbatim}

The check now expresses what it was always intended to express, namely whether
the Coder produced anything, and it does so correctly for both repair tasks and
generation tasks. The separate diff that is shown back to the Coder as a record
of its previous attempt is deliberately left as \texttt{git diff}, since that
value is genuinely required to be a diff.

\subsection{Context Between Coder Rounds}

The third intervention follows from the same assumption as the second, in a
different place.

Each Coder round begins a fresh conversation. What an earlier round did
reaches a later one only through a summary that the session assembles, built
from four parts: the report from the Tester, the diff of the previous attempt,
a record of the edits it made, and the file contents the Analyst read while
localising the work. On a repair task the diff alone carries the entire
attempt forward, and the arrangement is sound.

On a task built from scratch, three of those four parts are empty. The diff is
empty because nothing is tracked. The record of edits is empty because it is
derived from calls to the editing action, and a Coder writing whole files uses
the shell instead. The Analyst's file contents are empty because the Analyst
explored a workspace that did not yet contain anything. What remains is the
plan, the original specification, a test report, and an instruction stating
that the previous attempt was incorrect. A later round therefore rediscovers
its own work from scratch, while being told that work was wrong.

The effect is measurable. In a run of three rounds the Coder exhausted its
step budget in every one, consumed roughly two and a half times the input
tokens of an earlier single-round run of the same task, and produced strictly
less: the central module of the library was absent, and the benchmark score
fell from 0.505 to zero. Three restarts accomplished less than one
uninterrupted attempt. Other variables differed between those two runs, so
this is not a controlled comparison; the mechanism, however, is visible
directly in the assembled prompt.

We therefore give a later round an inventory of the workspace whenever the
diff is empty, listing the files that exist as
\texttt{\#\# Files already in the workspace}. The listing must be requested with
\texttt{--untracked-files=all}, since the default collapses an untracked
directory to its name and hides the files inside it, and it is capped so that
a single large generated directory cannot crowd out the rest of the prompt.
Where a diff exists the listing is suppressed and the prompt is assembled
exactly as before, so behaviour on repair tasks is unchanged.

This supplies information that the original design obtained from the diff. It
adds no instruction, grants no capability, and leaves the role prompts as
published.

\subsection{Sampling and Reasoning Parameters}

Every model call made by the Analyst and the Coder passes through a single
function; the Tester is deterministic and calls no model at all. That function
constructed its request from four values: the model, the messages,
\texttt{max\_tokens} and \texttt{temperature}. Two things follow from what it
omitted.

The first is that \texttt{top\_p} was dropped. The tool's own configuration
object declares it, and the sibling function used by the non-agentic runner
sends it; only the tool-calling path left it out. The same declared
configuration therefore produced different sampling depending on which entry
point was used, with the agentic path falling back to the provider's default.
We now send it on both paths. This does change behaviour relative to earlier
runs of the agentic path, which sampled at the provider's default rather than
at the value the configuration named, and we report it for that reason.

The second is that reasoning was never requested. For a model that exposes a
reasoning or thinking mode, the amount of reasoning performed was therefore
whatever the provider happened to default to for that model at that moment.
Such a default is set outside the experiment, does not appear in its record,
and may change between runs without notice. This is an uncontrolled variable in
a place where a comparison is particularly sensitive, because reasoning tokens
are drawn from the same output budget as the answer: a model that reasons at
length has correspondingly less budget left to act with, and the tool-calling
path has no fallback for a reply that consists of reasoning alone.

We therefore added two optional fields to the tool's model configuration and
pass them through to the request. \texttt{reasoning\_effort} carries the
OpenAI-compatible field of that name. \texttt{extra\_body} carries
provider-specific request fields verbatim, which accommodates interfaces that
express reasoning differently, such as OpenRouter's \texttt{reasoning} object.
Both are omitted from the request when unset, so a configuration that says
nothing about reasoning behaves exactly as before, and an unmodified revision
of the tool can still be pinned for a baseline run.

We then measured the defaults rather than assuming them. Both models used in
this work reason unprompted: on a trivial tool-calling request with nothing
asked for, \texttt{deepseek-v4-flash-0731} returned 16 reasoning tokens and
\texttt{kimi-k2.5} returned 40, in both cases accompanied by a reasoning field
on the message. Requesting a high effort, through either field, changed neither
model measurably on that request, while disabling reasoning through the
provider's own object removed it entirely. Every shape was accepted; none was
rejected. Our configurations accordingly pin reasoning explicitly to the state
the providers currently default to, so that the setting is fixed by the
experiment rather than by the provider, and an ablation without reasoning
becomes a one-line change. The probe was deliberately trivial, and establishes
which settings take effect rather than how they influence performance on a
real task.

Independently of whether reasoning is requested, each run now records the
number of reasoning tokens it consumed, taken from the usage the provider
reports. A run therefore carries evidence of what the model actually did, and
not only of what it was asked to do, which allows a run made under a provider
default to be identified as such after the fact. This matters in retrospect as
well as in prospect: runs made before this change consumed reasoning tokens
that were counted only in their output total, and cannot now be separated.

\subsection{Scope of the Intervention}

All three changes are confined to the harness that surrounds the roles: how it
inspects the repository, what it carries between rounds, and how it constructs
a request. The three role prompts, the tools available to each role, the
structure of the session loop and the order in which the roles run are
unmodified. No step was added to any role's instructions, and no role was given
a capability it did not previously have. The evaluated system is accordingly a
modified implementation rather than the published one, and we describe it as
such: the modifications adapt an implementation written for repair tasks to
tasks that begin from nothing, and none of them changes the method.

We nonetheless report them, rather than treating them as implementation
details, for two reasons. They are modifications to the system under
measurement, which should be disclosed on principle. And their practical effect
is not negligible. Without the first, the Tester role, though executed,
contributes no usable information, so a comparison drawn against the
unmodified implementation on this class of task would in effect be a
comparison against a two-role method. Without the second, the rounds after the
first are largely wasted, and additional rounds make the outcome worse rather
than better. Without the third, two runs of the same written configuration are
not guaranteed to have used the same sampling or the same amount of reasoning,
which is the property a controlled comparison depends on.

A pre-existing alternative deserves mention, since it was our first approach.
The same effect can be obtained without touching the implementation, by having
the Tester's command register the generated files with
\texttt{git add --intent-to-add} once the tests have run, which makes them
visible to \texttt{git diff}. We discarded it. It leaves the incorrect check in
place and compensates for it externally, and it has a side effect: once the
files are registered, the diff shown to the Coder in subsequent rounds contains
the entire generated project, inflating the prompt for no benefit. Correcting
the check is both smaller and more honest.

\section{Adapting CodeTeam}

The second solution under evaluation is CodeTeam, a multi-agent method in which
several \emph{Architect} agents each propose a software design sketch naming the
file tree, the public interfaces, the dependencies between files and the
developers the plan needs; a \emph{CTO} selects one of them and normalises it
into a contract; \emph{Developer} agents implement the files they own under a
dependency-aware scheduler, propagating interface changes to one another as
commits; and a \emph{QA} agent writes throwaway tests, runs them, and hands
failures back for repair until the repository converges or a budget is
exhausted. Unlike the other solution, the implementation we evaluate is the
authors' own repository at a pinned revision rather than a fork of it, for the
reason the next paragraph gives.

The account here is shorter than the preceding one, and for a reason worth
stating. The assumption that made three source changes necessary for
Self-Collaboration --- that a repository already exists, and that the work done
to it is visible as a diff --- does not arise here, because CodeTeam was written
for exactly this class of task: an empty workspace and a requirements document
from which a whole project must be produced. No file of CodeTeam is modified.
What this project supplies is an entry point and an instrument around it, and we
report those for the same reason we report the others.

\subsection{Where the Generated Repository Is Written}

CodeTeam's own entry point generates into a fresh directory named for the time
the run started, one level below its workspace. That suits a local run keeping
its attempts side by side. It does not suit this harness, because the directory
Harbor grades is the workspace itself: once the agent finishes, the task's
verifier copies what it finds at \texttt{/workspace} on top of the benchmark's
reference tests. Generated one level down, the entire project would be invisible
to the grader, which would score an empty repository.

The location is decided by a single method on the tool's context object, which
we override to return the workspace. Nothing else about the workflow changes.
The obvious alternative --- let the tool generate where it prefers and move the
result afterwards --- was rejected because the tool initialises a Git repository
at that root and commits to it as the developers work, so a move would either
carry that history to a location the tool never knew about or discard the record
of how the repository was assembled.

\subsection{Where the Tool's Own Record Is Written}

The same argument applies to a second directory. CodeTeam records what it did
--- every architect's candidate design, the CTO's choice and its rationale, the
normalised plan, and the result of each QA round --- beneath its workspace as
well. Under the arrangement above that places them inside the repository being
graded, where the verifier would copy them onto the benchmark's tests as though
they were part of the library under construction.

We direct them to the job's log directory instead, alongside the run's console
output. The alternative of switching the tool's artefacts off would remove the
problem and the evidence together, and the planning record is the part of this
method most worth having: it is the only place the designs that were rejected
can be read.

\subsection{Sampling and Reasoning Parameters}

The argument of the corresponding section above applies here too, and is
narrower in scope.

The inconsistency found in the other solution does not exist in this one. Every
request CodeTeam sends carries \texttt{temperature}, \texttt{top\_p} and
\texttt{max\_tokens} from one configuration object, on every path.

Two gaps remain. The first is that \texttt{temperature} has no environment
override, although the settings beside it do, so a run configured from outside
the tool could not fix its own sampling; this is among the reasons our runner
assembles the tool's configuration object directly rather than setting
environment variables and calling its entry point. The second is that no
reasoning field is ever sent. As before, the consequence is that the amount of
reasoning a run performs is whatever the provider happens to default to for that
model at that moment: a quantity set outside the experiment, absent from its
record, free to change between runs, and drawn from the same output budget as
the answer.

We add both by wrapping the single client method through which every one of the
tool's requests passes. \texttt{reasoning\_effort} carries the OpenAI-compatible
field of that name, and \texttt{extra\_body} carries provider-specific request
fields verbatim, which accommodates interfaces that express reasoning
differently. Both are omitted when unset, so a configuration that says nothing
about reasoning produces exactly the request the tool would have sent. The
tool's own request construction, its retries, and its JSON repair loop are
untouched.

\subsection{Usage and Rate Limits}

CodeTeam records one number about what it spends: a running total of total
tokens, which it consults only to enforce its own token budget. That is not
enough to report uncached input, cached input, output and reasoning tokens
separately, and those separations are what a comparison of cost between two
solutions rests on --- cached input is not priced as input, and reasoning tokens
are billed as output while forming no part of the answer.

The same wrapper therefore accumulates the usage each response reports, through
the same accounting the other solution uses. This is deliberately one
implementation rather than two: a token, a cached token and a dollar must mean
the same thing on both sides of a comparison, and two independent
implementations of that is how they quietly stop doing so.

The wrapper also waits out a throttled request. CodeTeam retries three times on
the path that returns free text and not at all on the path that returns JSON,
which is the majority of its calls, so a free model briefly out of capacity
would otherwise cost the trial. If every attempt is refused, the failure names
the rate limit and the model rather than surfacing as a JSON parsing error
several frames away.

\subsection{Budgets}

CodeTeam implements a wall-clock budget and a token budget and checks both
between stages. Neither is set by default. Left unset, the only bound on a task
is Harbor's own agent timeout, which is reached without anything having been
recorded about which stage was running or how much had been spent to get there.

Unlike the preceding subsections, this is a matter of configuration rather than
of implementation, as enabling Self-Collaboration's Tester was: we set both
values in the experiment configuration, where they are recorded with the run and
can be varied. It matters more here than it would for the other solution,
because CodeTeam has no fixed number of model calls. The width of the architect
search, the number of files the selected design names, and the number of QA
repair rounds each multiply what a task costs, and the second of those is chosen
by the model rather than by the configuration.

\subsection{Scope of the Intervention}

No file of CodeTeam is modified. The roles, their prompts, the order in which
they run, the scheduler, the Git-based coordination and the QA loop are the
published ones, and the revision evaluated is upstream's, pinned with the run
and recorded alongside its results. What this project adds surrounds the method
rather than entering it: where the finished repository is written, where the
tool's own record is written, which fields the request carries, what is counted,
and what bounds the run. The three components the paper ablates --- retrieval
grounding, dynamic developer allocation, and Git-based coordination --- are
exposed as settings of the experiment configuration, each recorded with the run
it produced.

\section{Adapting CodeS}

The third solution under evaluation is CodeS, which is not a team but a
pipeline. It decomposes the construction of a repository into three layers of
sketch: a \emph{RepoSketcher} proposes the file tree implied by the
requirements, a \emph{FileSketcher} writes each Python file that tree names as
signatures with empty bodies, and a \emph{SketchFiller} implements one function
per request, given that file's sketch and the sketches of the files it imports.
A final stage parses each sketch, substitutes the generated bodies into it, and
writes the result to disk. No agent talks to another and nothing is revised: a
request is made once, in one of three shapes, and the pipeline runs forward.

As with CodeTeam, the implementation we evaluate is the authors' own repository
at a pinned revision rather than a fork, and no file of CodeS is modified. The
account below is longer than CodeTeam's nonetheless, and for a reason that has
nothing to do with the quality of the implementation. CodeTeam ships a program;
CodeS ships a research artefact, in which the pipeline is expressed as a script
rather than as a library, and in which the method as published is a fine-tuned
model as much as it is the framework around it. Both facts bear on what a
result obtained here can be said to be about, and the first two subsections are
therefore about what is being measured rather than about how.

\subsection{Which CodeS Is Evaluated}

CodeS is published as a fine-tuned model together with the framework that
prompts it. The paper's contribution is both: the multi-layer sketch is the
decomposition, and the released checkpoints are models trained on instruction
data extracted from a hundred public repositories to answer each of the three
layers in turn.

The tool accordingly ships two drivers for the same pipeline. Its principal one
loads the fine-tuned checkpoint locally through \texttt{transformers} and
generates with a text-generation pipeline on a GPU; it reaches no API of any
kind. Its second, distributed among the baselines, runs the identical three
phases against an OpenAI-compatible endpoint, calling the same phase helpers
and the same assembly stage, and differing only in the prompt wording that a
chat model rather than a completion-trained one requires.

We integrate the second. The reason is the platform's, not the method's: this
harness supplies one credential and one base URL, and a comparison across
solutions rests on all of them reaching the same provider in the same way. A
locally hosted checkpoint would be a second model backend in an experiment
whose entire design is that the model is held constant.

The consequence must be stated plainly, because it is the largest single
qualification on any result reported for this solution. What is measured here
is CodeS's multi-layer sketch driven by the same general-purpose model as the
other two solutions --- which is a configuration the paper itself reports ---
and not the fine-tuned CodeS model. A comparison drawn here is therefore
between three frameworks at one model, and it is silent on the contribution of
the training that gives CodeS its name. We regard this as the right comparison
for this work, since the other two solutions are prompting frameworks over a
general-purpose model and nothing else would place them on the same axis, but
it is a narrower claim than ``CodeS was evaluated'' and should not be reported
as the wider one.

\subsection{Driving a Script}

The driver we integrate is a script and not a library. Its prompt templates,
its argument parsing, its three-phase loop and its output paths are all
statements at module scope, so that importing it parses a command line, opens
files relative to the working directory, and begins making requests. There is
no function to call.

Our runner therefore repeats the loop and takes everything else from the tool.
The decomposition is not reimplemented: which paths the repository sketch
names, which of them are sketched, which functions are asked for, what context
each prompt carries, how a response is parsed, and how the answers are
assembled into files are all calls into CodeS's own modules, unchanged. What
the runner supplies is the sequencing between those calls, which is the part
that existed only as top-level statements.

The prompts are the one thing that could not be reached by a call, and copying
them into this repository was rejected. A copy would be free to disagree with
the revision a configuration pins, and nothing in the record would say which of
the two a run had used --- precisely the failure that pinning a revision exists
to prevent. The runner instead reads the driver's source, locates the literal
assignment that holds the templates, and evaluates that assignment alone
without executing the module. The prompts a run used are thus the pinned
revision's by construction, and a revision that no longer carries them fails
immediately and says so rather than falling back to a stale copy.

\subsection{Where the Generated Repository Is Written}

The argument is CodeTeam's, and the mechanism differs. CodeS's assembly stage
writes each repository into a directory named after it, one level below the
output directory it is given, which suits its own evaluation --- a hundred
generated repositories side by side for scoring. Harbor grades
\texttt{/workspace} itself, so a project written one level down would be scored
as an empty repository rather than reported as misplaced.

Unlike CodeTeam, whose generation root is decided by a single method we could
override, CodeS decides it in the loop of its assembly stage, which is again
top-level script code. Our runner therefore calls the three functions that
stage is built from --- the one that reads the file sketches, the one that
groups the generated bodies by file, and the one that substitutes the bodies
into a sketch --- and writes the results itself. The substitution logic, which
is the part that matters, is the tool's and is untouched.

Two safeguards are added at the point of writing. The paths come from a file
tree the model wrote, so a path that resolves outside the workspace is dropped
and named in the run's record rather than written where it points; and each
file is assembled independently, for the reason a later subsection gives.

\subsection{Where the Tool's Own Record Is Written}

CodeS records each phase as a file of prompts and answers: the repository
sketch it proposed, the sketch of every file, and every function it filled in.
It writes them into the directory it generates into, which under the
arrangement above would place them inside the repository being graded, where
the verifier would copy them onto the benchmark's reference tests. They are
directed to the job's log directory instead, alongside the run's console
output.

The record is more valuable for this solution than for the others, and is worth
keeping rather than merely relocating. CodeS never revises: a function is asked
for once and the answer is used. The prompt that produced a given function is
thus the complete explanation of it, and the assembled repository is a
deterministic function of these files.

\subsection{Sampling and Reasoning Parameters}

The tool's driver sends a request carrying the model, the messages, and a
temperature, and nothing else. It gives each request five attempts, the first
at temperature $0.0$ and the remainder at $0.1$, and calls \texttt{exit(1)}
when all five are refused. Every one of those figures is a literal in the
script.

Three things follow. The sampling of a run could not be part of its recorded
setup, since neither the temperatures nor the attempt count is reachable from
outside; they are hyperparameters here, defaulting to the tool's own values, so
that a configuration stating nothing reproduces the published behaviour and one
stating something records what it changed. Second, \texttt{top\_p} is never
sent, so that dimension of sampling is the provider's default: a quantity set
outside the experiment and absent from its record. Third, no reasoning field is
ever sent, with the consequence described at length for the other two
solutions.

A fourth omission is specific to this method. The driver sends no token limit,
so the length of a response is bounded by whatever the provider defaults to.
For a solution that returns a patch this is a minor risk; for CodeS it is not,
because a response is a whole file sketch or a whole function. A response
truncated mid-definition does not produce a weaker implementation --- it
produces text that does not parse, and the file it belonged to is written as an
unfilled sketch instead. Our configurations set a limit for that reason.

All four are supplied by constructing the request the tool constructs and
adding the fields it omits, each omitted again when unset, so that a
configuration that states nothing about them sends exactly the request the tool
would have sent.

\subsection{Usage and Rate Limits}

CodeS records nothing at all about what it spends. The usage accounting is
therefore entirely supplied, through the same shared implementation the other
two solutions use, for the reason given there: a token, a cached token and a
dollar must mean the same thing on both sides of a comparison.

The tool's own retry loop needed one change. It retries any failed request,
which already covers a throttled one, but it retries immediately --- so a
provider briefly out of capacity is met with five requests in a few seconds and
the trial is lost. We insert a wait before an attempt that follows a throttled
one, and deliberately inside the tool's attempt count rather than beside it, so
that a rate limit costs time and not a multiple of the requests the
configuration asked for. A failure that is not a rate limit is retried
immediately, as before, since a rejected credential or an unknown model will be
refused just as quickly the second time.

One further condition is treated as a failed attempt rather than as an answer.
A response whose content is empty --- filtered, truncated before any text, or
consisting of reasoning with nothing after it --- would otherwise be parsed
into an empty file sketch or an empty function body and carried forward as
though the model had answered. Retrying costs one request; not retrying costs
the file.

\subsection{Budgets, and What a Stopped Run Leaves Behind}

CodeS has no bound of its own on what a task costs, and unlike the other two
solutions it cannot be bounded by asking for fewer rounds. Its length is a
property of the design the model proposed in its first response: one request,
plus one for each Python file the repository sketch named, plus one for every
function those file sketches declared. A specification that invites a wide
design costs several times what a narrow one does. As with CodeTeam we set a
wall-clock budget and a token budget in the experiment configuration, where
they are recorded with the run and can be varied.

The budget interacts with a structural property of this method that required a
change. CodeTeam builds the repository as it goes, so a run stopped part-way
has already written files. CodeS assembles only at the very end, so a run
stopped part-way has written nothing --- and its assembly stage, given a record
of a request that was never answered, fails on that record and produces no
repository at all. The whole of a run's expenditure would be lost to the
requests it did not make.

The runner therefore always assembles, whether the pipeline finished, exhausted
its budget, or failed outright, and passes the assembly stage only the requests
that were answered. This is less a change of behaviour than a restoration of
one: the tool appends each answer to these files as it arrives, so an
interrupted run of the tool leaves exactly the files we construct, and our
writing them once at the end is what introduced the discrepancy. What a stopped
run leaves is a repository of correct interfaces with some bodies left as
\texttt{pass}, which the verifier can grade, in place of an empty workspace
which it can only score as nothing.

\subsection{Containing a Failure to the File It Came From}

Two places in the tool process every file in a single unguarded loop, so that
one file ends the loop and takes the rest with it.

The first is the construction of the third phase's prompts, which parses each
file sketch to find the functions it declares. The tool tolerates a sketch that
is merely truncated, trimming lines until what remains parses, but a sketch
malformed in other ways raises out of the loop and the remaining files never
have their function bodies requested. The second is the assembly stage, where a
sketch that defeats the substitution ends the assembly with the repository
partly written.

Each is enclosed so that the failure costs the file it came from and no other:
in the first case that file's functions are not requested, and in the second
the file is written as the sketch FileSketcher produced rather than not written
at all. Both counts are recorded with the run, so a repository that is thin
because of this is distinguishable after the fact from one that is thin because
the model proposed little. We report these as modifications, though they add no
capability and change no output the tool would itself have produced, because
their effect on a measured score is not small: under the unmodified control
flow a single malformed response can reduce a run to an empty repository.

\subsection{Concurrency Within a Phase}

The published pipeline is strictly sequential, and a repository of ten modules
declaring sixty functions is seventy-one requests. In sequence that exceeds the
agent timeout of the tasks used here before the pipeline can finish, which is
scored as an error rather than as a poor repository --- a measurement of the
harness rather than of the method.

We therefore allow the requests within one phase to be in flight together, as a
configuration setting whose default is $1$, which is the published pipeline
exactly. This is admissible because of what the phases are, and would not be
admissible for either other solution: every file sketch reads the one
repository sketch and nothing else, and every function body reads the completed
set of file sketches, which the tool itself collects before the third phase
begins. No request in a phase can observe another request in the same phase
under any ordering, so the set of prompts sent and the mapping from prompts to
answers are unchanged; what changes is how long the phase takes. The records
are kept in the order the sequential driver would have produced them.

The value used is recorded with each run and reported with our results, since
it is a deviation from the published pipeline even though it is one we can
argue is inert.

\subsection{The Runtime the Tool Requires}

Two properties of the checkout are worth recording, since both bear on
reproducing a run rather than on the method.

CodeS builds and inspects syntax trees using \texttt{ast.Str}, an alias Python
deprecated in version 3.8 and removed in 3.14. Where it is absent the tool
raises from inside the prompt construction of its third phase, after the
repository sketch and every file sketch have already been paid for. The runner
checks for it before making any request and fails with that explanation
instead, and the task image is accordingly required to carry Python 3.13 or
older for this solution.

The repository is also large for what a run reads: it carries the hundred
public repositories its training data was extracted from and the hundred more
its benchmark grades, none of which any run opens, and every trial clones it.
The clone is made partial and sparse, fetching the commits and trees but file
contents only for the directory the pipeline reads. This narrows the working
tree and not the revision, which is pinned exactly as before, and reduces a
per-trial cost that a benchmark run pays several hundred times.

\subsection{Scope of the Intervention}

No file of CodeS is modified. The three prompts are read from the pinned
revision rather than reproduced, and the decomposition, the parsing of each
response, the selection of context for each prompt and the substitution of
bodies into sketches are the published implementation's, called rather than
reimplemented. What this project adds surrounds the method: the sequencing that
existed only as top-level script statements, where the finished repository and
the tool's own record are written, which fields the request carries, what is
counted, what bounds the run, that a stopped run still assembles, that one
file's failure is not every file's, and that a phase's requests may overlap.

The qualification that matters most is not in that list, and is stated again
here because it is easily lost among the others: the model driven by this
framework is the experiment's model and not the fine-tuned model the paper
released. A result reported for this solution is a result for the multi-layer
sketch, at the model held constant across the comparison.
